\documentclass[11pt,letterpaper]{article}

\usepackage[spanish,es-noshorthands]{babel}
\usepackage{fontspec}
\setmainfont{Source Serif Pro}
\usepackage{microtype}
\usepackage{amsmath}
\usepackage{geometry}
\usepackage{booktabs}
\usepackage{tabularx}
\usepackage{enumitem}
\usepackage{xcolor}
\usepackage{graphicx}
\usepackage{csquotes}
\usepackage{tikz}
\usetikzlibrary{arrows.meta,positioning,shapes.geometric,calc}
\usepackage{pgfplots}
\pgfplotsset{compat=1.18}
\usepackage{xurl}
\usepackage[hidelinks]{hyperref}
\usepackage[backend=biber,style=apa,sorting=nyt]{biblatex}
\DeclareLanguageMapping{spanish}{spanish-apa}

\geometry{margin=2.5cm}
\setlength{\emergencystretch}{1em}
\setlength{\parindent}{0pt}
\setlength{\parskip}{0.65em}
\setlist{nosep}
\addbibresource{referencias.bib}
\AtBeginBibliography{\sloppy}

\title{\textbf{Comentario sobre una hoja de ruta tecnológica}\\
\large \emph{Quantum Machine Learning}: cobertura de los doce elementos y
viabilidad de proyectar una fecha de disponibilidad}
\author{Carlos Luis Rojas Aragonés\\
Gestión del Desarrollo Tecnológico}
\date{27 de agosto de 2026}

\begin{document}

\maketitle

\section{Objeto, fuente y método del comentario}

\subsection*{01. El documento analizado}

Este comentario examina la hoja de ruta tecnológica titulada \emph{Quantum
Machine Learning}, elaborada por Kentaro Numa, Dian Wen y Jessy M. Mwarage para
el curso \emph{16.887/EM.427 Technology Roadmapping and Development} en el
otoño de 2023 \parencite{numa2023}. El archivo es una exportación de
PowerPoint de nueve páginas, con el título interno \enquote{Quantum ML -
Technology Overview\_2} y fecha de creación del 5 de diciembre de 2023. De esas
nueve páginas, una es la portada y ocho contienen material analítico.

La hoja de ruta plantea el caso de una empresa que pretende lanzar un producto
\emph{SaaS} de aprendizaje automático cuántico para clientes B2B, combinando
hardware cuántico adquirido y desarrollado internamente. Su declaración de
estrategia fija el objetivo en el año 2035.

\subsection*{02. Método seguido}

El comentario se construyó en tres pasadas sucesivas:

\begin{enumerate}
  \item \textbf{Transcripción completa.} Se extrajo la capa de texto del
    archivo y se leyeron a alta resolución todas las figuras que en la
    presentación aparecen reducidas: el diagrama OPM, la matriz DSM, la lista
    competitiva de veinte procesadores, la matriz morfológica, el análisis de
    sensibilidad y las cinco tablas del modelo financiero. Sin esa lectura
    ampliada, buena parte del contenido cuantitativo de la hoja de ruta
    resulta ilegible.
  \item \textbf{Reproducción aritmética.} Se recalcularon el modelo de volumen
    cuántico, las derivadas tecnológicas normalizadas y la totalidad del flujo
    de caja descontado, para comprobar si los resultados publicados se
    sostienen a partir de los datos y supuestos que la propia hoja de ruta
    declara.
  \item \textbf{Corroboración externa.} Toda afirmación que no proviene del
    archivo se verificó contra la fuente primaria. Esto incluye las
    referencias citadas por la hoja de ruta, cuyos metadatos se comprobaron en
    Crossref y en los registros de patentes.
\end{enumerate}

A lo largo del texto se distingue de forma explícita entre lo que la hoja de
ruta afirma, lo que se ha podido verificar y lo que constituye una derivación
propia de este comentario a partir de los datos contenidos en ella.

\section{Qué contiene la hoja de ruta}

El Cuadro~\ref{tab:contenido} resume el contenido de las ocho páginas
analíticas.

\begin{table}[ht]
  \centering
  \small
  \caption{Contenido de las páginas analíticas de la hoja de ruta.}
  \label{tab:contenido}
  \begin{tabularx}{\textwidth}{@{}c>{\raggedright\arraybackslash}X@{}}
    \toprule
    \textbf{Pág.} & \textbf{Contenido} \\
    \midrule
    2 & \emph{Roadmap Overview}. Clasificación tecnológica
        (\enquote{Process Information}), descripción de la computación cuántica
        y del cúbit, esquema comparativo entre aprendizaje automático clásico y
        cuántico, diagrama OPM y matriz DSM. \\
    3 & \emph{FOMs and Technology Evolution}. Lista de nueve figuras de mérito
        y tres gráficos tomados de fuentes externas. \\
    4 & \emph{Technology Strategy Statement}. Objetivo a 2035 y tres líneas
        estratégicas con hitos fechados. \\
    5 & \emph{Key Publications \& Patents}. Genealogía de publicaciones entre
        1995 y 2021 y dos patentes. \\
    6 & \emph{Alignment with Company Drivers}. Pila de sistema en cuatro capas
        y tabla de tres motivaciones estratégicas con cuatro objetivos
        cuantificados y semáforo de alineación. \\
    7 & \emph{R\&D Projects, Company Positioning vs. Competition}. Mapa de
        correlación generado con inteligencia artificial y lista competitiva
        de veinte procesadores cuánticos. \\
    8 & \emph{Technical Model}. Matriz morfológica de tres opciones de
        hardware y análisis de sensibilidad del volumen cuántico. \\
    9 & \emph{Financial Model}. Análisis de compañías comparables, proyección
        anual 2022--2035 y valor presente neto. \\
    \bottomrule
  \end{tabularx}
\end{table}

\section{Cobertura de los doce elementos}

La estructura recomendada para una hoja de ruta tecnológica contempla doce
elementos. El Cuadro~\ref{tab:cobertura} indica dónde se localiza cada uno y
qué valoración merece.

\begin{table}[ht]
  \centering
  \footnotesize
  \caption{Cobertura de los doce elementos en la hoja de ruta analizada.}
  \label{tab:cobertura}
  \begin{tabularx}{\textwidth}{@{}c>{\raggedright\arraybackslash}p{0.26\textwidth}
    c>{\raggedright\arraybackslash}X@{}}
    \toprule
    \textbf{N.\textsuperscript{o}} & \textbf{Elemento} & \textbf{Pág.} &
    \textbf{Valoración} \\
    \midrule
    1 & Descripción de la hoja de ruta & 2 & Presente. Clara en la
      descripción física, ausente en la delimitación del alcance temporal. \\
    2 & Matriz de estructura de dependencia (DSM) & 2 & Presente, pero
      cataloga componentes internos y no otras hojas de ruta, que es lo que su
      propio título anuncia. \\
    3 & Modelo de la hoja de ruta (OPM) & 2 & Presente. Diagrama de objetos y
      procesos válido en su topología, con dos relaciones de especialización
      invertidas y sin atributos. \\
    4 & Figuras de mérito (FOM) & 3 & Presente sólo como nómina. Ninguna de
      las nueve tiene valor de partida, valor objetivo ni serie temporal. \\
    5 & Correspondencia con las motivaciones estratégicas & 6 & Presente y es
      el elemento mejor resuelto del conjunto. \\
    6 & Posicionamiento competitivo mediante diagramas de FOM & 7 & Parcial.
      Existe la tabla competitiva, pero no hay ningún diagrama de figuras de
      mérito. \\
    7 & Modelo técnico & 8 & Presente. Contiene un error de base logarítmica,
      una inconsistencia de unidades y un punto de diseño fuera del espacio
      que él mismo define. \\
    8 & Modelo financiero & 9 & Presente y aritméticamente reproducible, con
      supuestos de entrada que presuponen la respuesta. \\
    9 & Portfolio de proyectos y prototipos de I+D & 7 & \textbf{Ausente.} La
      página que lo anuncia muestra hardware de la competencia, no proyectos
      propios. \\
    10 & Publicaciones, presentaciones y patentes clave & 5 & Presente.
      Genealogía útil, sin presentaciones y con una fecha de patente
      incorrecta. \\
    11 & Declaración de estrategia tecnológica & 4 & Presente. Es el único
      lugar donde aparecen fechas objetivo. \\
    12 & Evaluación del nivel de madurez de la hoja de ruta & n/d &
      \textbf{Ausente.} \\
    \bottomrule
  \end{tabularx}
\end{table}

En balance, nueve elementos están presentes de forma sustantiva, uno está
resuelto a medias y dos faltan por completo. La observación relevante para el
objeto de este comentario es que los dos elementos ausentes, el portfolio de
proyectos y la evaluación de madurez, son precisamente los dos que contienen
información de calendario y de disposición a producir. Su ausencia no es una
laguna formal: es la razón principal por la que la hoja de ruta no puede
sostener una proyección de disponibilidad.

\section{Los elementos de modelado (1 a 3)}

\subsection*{01. Descripción y clasificación}

La descripción cumple su función. Sitúa el aprendizaje automático cuántico como
la intersección entre la computación cuántica y el aprendizaje automático,
explica la superposición y advierte que distintas realizaciones físicas del
cúbit, superconductor u óptico, condicionan la teoría y el aparato de cálculo.
Clasifica la tecnología como \enquote{Process Information}, es decir, en
función de la operación genérica que ejecuta sobre la materia, la energía o la
información.

La carencia es de otro tipo. La descripción no delimita qué se entiende por
disponibilidad de la tecnología. No dice si el objeto de la hoja de ruta es un
acelerador cuántico para cargas de trabajo concretas, una biblioteca de
algoritmos, un servicio en la nube o un producto integrado. Sin esa
delimitación no hay un umbral que fechar, y esta indefinición reaparece en cada
elemento posterior.

Tampoco se declara el identificador ni el nivel de la propia hoja de ruta
dentro de una jerarquía. Es un detalle revelador, porque la página 6 sí utiliza
el identificador \enquote{1QCAIML} para la hoja de ruta de nivel superior al
encabezar la columna de objetivos. El vocabulario jerárquico existe en el
documento, pero no se aplica al documento mismo.

\subsection*{02. El modelo OPM}

El diagrama de objetos y procesos es correcto en su topología general.
Distingue once objetos, dibujados como rectángulos, y cinco procesos,
dibujados como elipses. Los objetos son cúbits, puertas cuánticas, circuitos
cuánticos, entrelazamiento cuántico, algoritmos cuánticos, VQA, datos
cuánticos, LDA, modelos de aprendizaje automático, QNN y resultados. Los
procesos son \emph{Computing}, \emph{Encoding}, \emph{Training},
\emph{Predicting} y \emph{Analyzing}. La cadena funcional que describe es
legible: los recursos cuánticos habilitan el cómputo, el cómputo actúa sobre
los algoritmos, los algoritmos entrenan modelos y los modelos predicen y
analizan hasta producir resultados.

Hay tres observaciones que conviene registrar.

Primera, las dos relaciones de especialización están invertidas. El triángulo
que vincula VQA con los algoritmos cuánticos tiene su vértice apuntando hacia
VQA, y lo mismo ocurre entre QNN y los modelos de aprendizaje automático. Leído
según la convención de la notación, el diagrama afirma que un algoritmo
cuántico es un tipo de VQA y que un modelo de aprendizaje automático es un tipo
de QNN, cuando la taxonomía pretendida es la contraria.

Segunda, el tratamiento de las familias de algoritmos es incoherente. VQA y QNN
se enganchan como especializaciones, mientras LDA, que también es una familia
de métodos, aparece como operando del proceso \emph{Encoding}. Tres objetos del
mismo tipo lógico reciben tres tratamientos distintos.

Tercera, y con mucho la más importante, ningún objeto del diagrama tiene
atributos. En un modelo de objetos y procesos, los atributos son el lugar
natural donde se anclan las figuras de mérito: la fidelidad es un atributo de
las puertas cuánticas, el número de cúbits lo es de los cúbits, la exactitud lo
es de los modelos. Al no declarar atributos, el diagrama describe qué hace el
sistema pero no qué se mide en él, y queda desconectado de las nueve figuras de
mérito enumeradas en la página siguiente. Tampoco se acompaña de la
descripción textual equivalente que la metodología prevé, ni de niveles de
detalle sucesivos.

Se aprecia además una errata en el diagrama, \enquote{Quantum Engtanglement}.

\subsection*{03. La matriz DSM}

La matriz se titula \enquote{DSM Allocation (interdependencies with others
roadmaps)}. El título anuncia interdependencias con otras hojas de ruta, pero
sus once filas y once columnas son exactamente los once objetos del diagrama
OPM. Lo que la matriz representa son dependencias entre componentes internos
del sistema, no entre hojas de ruta. Ninguno de los cinco procesos del
diagrama aparece en ella, y tampoco aparece el identificador \enquote{1QCAIML}
que la página 6 sí emplea. La consecuencia práctica es que la hoja de ruta no
documenta de qué otras hojas de ruta depende, y por tanto no documenta qué
esperas ajenas condicionan su propio calendario.

La reproducción de la matriz permite tres constataciones adicionales. Contiene
44 marcas sobre 110 celdas fuera de la diagonal, una densidad del 40 por
ciento. De los 22 pares posibles con marca, 20 son simétricos y 4 no lo son:
resultados frente a VQA, resultados frente a QNN, resultados frente a LDA y
datos cuánticos frente a LDA. Las cuatro asimetrías apuntan en el mismo
sentido, de un componente hacia los resultados, lo que sugiere que se intentó
registrar una dirección. Pero si la matriz fuera dirigida, los otros veinte
pares simétricos serían incorrectos, y si fuera no dirigida, esas cuatro
celdas están incompletas. La convención no se declara y se aplica de las dos
maneras a la vez.

Por último, las marcas son binarias. No hay pesos, ni tipos de dependencia, ni
orden de precedencia. Una DSM binaria permite ver acoplamientos, pero no
permite agrupar módulos por intensidad ni derivar una secuencia de desarrollo,
que es el uso que la convertiría en un instrumento de planificación temporal.

\section{El problema central: tres vocabularios de métricas desconectados}

Este es, a juicio de quien comenta, el defecto estructural del documento y la
causa directa de que no pueda proyectarse una fecha. La hoja de ruta utiliza
tres conjuntos de métricas que no se solapan entre sí.

\subsection*{01. Las nueve figuras de mérito declaradas}

La página 3 enumera nueve figuras de mérito con sus unidades: aceleración
computacional en FLOPS, eficiencia de aprendizaje medida como valor de pérdida
en por ciento, exactitud del modelo en por ciento, escalabilidad en datos por
segundo, versatilidad como número, generalización en por ciento, eficiencia de
recursos expresada en \enquote{valor}, CAPEX en dólares y OPEX en dólares.

Ninguna de las nueve tiene valor actual, valor objetivo, método de medición ni
serie histórica. Son nombres con unidades. Además, algunas unidades no
corresponden a la magnitud: una aceleración es un cociente adimensional y no se
mide en FLOPS, que son un caudal de operaciones; un valor de pérdida no es en
general un porcentaje; y \enquote{valor} no es una unidad. La escalabilidad
expresada en datos por segundo se solapa conceptualmente con la aceleración
computacional.

\subsection*{02. Los objetivos estratégicos}

La página 6 traduce tres motivaciones de la empresa en cuatro objetivos
cuantificados, y les asigna un semáforo de alineación: más de 16 cúbits
(alineado), fidelidad de cúbit superior al 99,9 por ciento (alineado), más del
1 por ciento de cuota de mercado del negocio \emph{SaaS} de aprendizaje
automático (en riesgo) y menos del 25 por ciento de variabilidad en el coste de
adquisición de hardware (no alineado).

Este es el elemento mejor construido de la hoja de ruta. Los objetivos son
verificables, el semáforo es honesto y reconoce dos incumplimientos, y el texto
razona por qué la fidelidad importa más que el recuento de cúbits. Sin embargo,
ninguno de los cuatro objetivos figura entre las nueve figuras de mérito de la
página anterior.

Hay además una imprecisión con consecuencias. El objetivo de fidelidad no
especifica si se refiere a puertas de un cúbit o de dos. El dispositivo elegido,
el H1-1 de Quantinuum, declara en la propia tabla de la hoja de ruta un 99,996
por ciento para un cúbit y un 99,8 por ciento para dos cúbits. Leído sobre
puertas de un cúbit, el objetivo se cumple y el semáforo verde es correcto;
leído sobre puertas de dos cúbits, no se cumple y el semáforo debería ser rojo.
Como el propio diagrama OPM incluye el entrelazamiento cuántico entre los
objetos necesarios para el cómputo, y el entrelazamiento se genera con puertas
de dos cúbits, la lectura pertinente es la segunda. La ambigüedad de una sola
palabra decide el color de la casilla.

\subsection*{03. El volumen cuántico}

Las páginas 7 y 8 introducen un tercer vocabulario. El texto afirma que el
volumen cuántico, entendido como un equilibrio entre número de cúbits y tasa de
error, es \enquote{la métrica principal para evaluar el rendimiento de un
computador cuántico}, y sobre él se decide la selección del dispositivo. El
volumen cuántico no aparece entre las nueve figuras de mérito. Tampoco la
fidelidad ni el número de cúbits.

La Figura~\ref{fig:vocabularios} representa la situación. Tres conjuntos de
métricas, cada uno definido en una página distinta, sin ninguna relación
declarada entre ellos. La estrategia se evalúa con unas métricas, el sistema se
describe con otras y la decisión técnica se toma con unas terceras. Falta la
trazabilidad que permitiría afirmar que alcanzar cierto volumen cuántico
implica alcanzar cierta exactitud de modelo, y que esa exactitud es la que
satisface la motivación estratégica. Sin esa cadena, no existe un umbral
técnico asociado a la disponibilidad comercial, y por tanto no hay nada que
fechar.

\begin{figure}[ht]
  \centering
  \begin{tikzpicture}[
    caja/.style={draw, rounded corners, align=left, text width=0.28\textwidth,
      inner sep=5pt, font=\footnotesize},
    roto/.style={thick, dashed, gray}
  ]
    \node[caja, fill=blue!6] (fom) {\textbf{Pág. 3: nueve FOM}\\
      Aceleración [FLOPS]\\ Eficiencia de aprendizaje [\%]\\
      Exactitud [\%] · Escalabilidad\\ Versatilidad · Generalización\\
      Eficiencia de recursos\\ CAPEX · OPEX};
    \node[caja, fill=orange!10, anchor=north west]
      at ([xshift=0.5cm]fom.north east) (drv)
      {\textbf{Pág. 6: objetivos}\\ $>16$ cúbits\\ Fidelidad $>99{,}9\%$\\
       Cuota de mercado $>1\%$\\ Variabilidad de coste $<25\%$};
    \node[caja, fill=green!9, anchor=north west]
      at ([xshift=0.5cm]drv.north east) (qv)
      {\textbf{Págs. 7--8: modelo}\\ Volumen cuántico (QV)\\
       Número de cúbits $n$\\ Tasa de error $p$};
    \draw[roto] ($(fom.north east)+(0,-9mm)$) --
      node[midway, fill=white, inner sep=1pt, font=\footnotesize,
        text=red!65!black] {$\times$} ($(drv.north west)+(0,-9mm)$);
    \draw[roto] ($(drv.north east)+(0,-9mm)$) --
      node[midway, fill=white, inner sep=1pt, font=\footnotesize,
        text=red!65!black] {$\times$} ($(qv.north west)+(0,-9mm)$);
  \end{tikzpicture}
  \caption{Los tres vocabularios de métricas de la hoja de ruta. Ningún
  elemento de un conjunto reaparece en otro y la hoja de ruta no declara
  ninguna relación entre ellos, lo que impide trazar un objetivo estratégico
  hasta la magnitud técnica que lo mediría. Elaboración propia a partir de
  \textcite{numa2023}.}
  \label{fig:vocabularios}
\end{figure}

\subsection*{04. El posicionamiento competitivo}

El elemento 6 pide diagramas de figuras de mérito para situar a la empresa
frente a la competencia. La hoja de ruta aporta una tabla de veinte
procesadores con fabricante, fidelidad, cúbits físicos, volumen cuántico y año
de lanzamiento, pero ningún diagrama. Es una omisión llamativa porque los datos
para construirlo están en la misma página: bastaría representar cúbits frente a
tasa de error, o volumen cuántico frente a año, para obtener la frontera de
diseño y la posición relativa del dispositivo elegido.

La tabla tiene además dos limitaciones de origen. La hoja de ruta declara que
proviene de la lista de procesadores cuánticos de Wikipedia, y el extracto
reproducido empieza en IonQ y termina en el Aspen-M-3 de Rigetti. Es un
fragmento alfabético: faltan IBM, Google, Intel y D-Wave. La ausencia de IBM y
Google es particularmente incómoda porque son dos de las cuatro compañías que
el modelo financiero de la página 9 utiliza como comparables. La segunda
limitación es que la columna de volumen cuántico sólo está rellena en cuatro de
las veinte filas.

El acompañamiento visual de esa página es un \enquote{Correlation Map of
Quantum Computer R\&D Projects} que, según la nota al pie de la propia hoja de
ruta, se generó con DALL-E 2 a partir de la lista de Wikipedia. Una imagen
generativa no codifica estructura de correlación alguna: es una ilustración
ocupando el lugar de un análisis.

\section{El modelo técnico (elemento 7)}

\subsection*{01. La matriz morfológica}

La matriz compara tres dispositivos de Quantinuum sobre tres variables: número
de cúbits $n$, tasa de error $p$ en por ciento y volumen cuántico. Los valores
son H2 con 32 cúbits y $2^{16}$, H1-1 con 20 cúbits y $2^{19}$, y H1-2 con 12
cúbits y $2^{12}$. Estos valores coinciden con la tabla competitiva de la
página anterior, y los rangos de $p$ se derivan correctamente de las
fidelidades declaradas. Hasta aquí, el modelo es consistente.

La decisión que se toma sobre esa matriz, en cambio, no se sigue de los
criterios que la hoja de ruta había fijado. Se elige el H1-1 por tener el
volumen cuántico más alto. Pero los dos objetivos que la página 6 establece
para el hardware son más de 16 cúbits y fidelidad superior al 99,9 por ciento,
y sobre esos dos criterios el H2 domina al H1-1: tiene 32 cúbits frente a 20,
igual fidelidad de dos cúbits (99,8 por ciento), mejor fidelidad de un cúbit
(99,997 frente a 99,996 por ciento) y es un año más reciente. El criterio que
decide, el volumen cuántico, se introduce en el momento de decidir y no figura
ni entre los objetivos estratégicos ni entre las figuras de mérito.

\subsection*{02. El análisis de sensibilidad y la base del logaritmo}

El análisis de sensibilidad calcula derivadas tecnológicas normalizadas del
volumen cuántico respecto a $n$ y a $p$, sobre un modelo que, reconstruido a
partir de las expresiones y de la tabla numérica, es
\[
  QV = 2^{n}\,(1-p)^{2}.
\]
La derivada respecto a $p$ está bien calculada. La hoja de ruta obtiene
$\partial QV/\partial p = -2^{\,n+1}(1-p)$ y, normalizada, $-2/(1-p)$, que en el
punto nominal $p = 5$ por ciento vale $-2{,}105$. Reproduce sus propias cifras.

La derivada respecto a $n$ contiene un error de base logarítmica. La hoja de
ruta emplea $\partial QV/\partial n = 2^{n}\log(2)\,(1-p)^{2}$ y normaliza a
$\log(2) = 0{,}301$, es decir, logaritmo decimal. La derivada de $2^{n}$
respecto a $n$ es $2^{n}\ln 2$, de modo que el valor normalizado correcto es
$\ln 2 = 0{,}693$. La discrepancia es un factor de $\ln 2 / \log_{10} 2 =
2{,}303$.

El error no invierte la conclusión, pero sí la magnifica. Con las cifras
publicadas, la tasa de error parece tener casi siete veces más apalancamiento
que el número de cúbits ($2{,}105 / 0{,}301 = 6{,}99$). Corregida la base, el
factor es de tres ($2{,}105 / 0{,}693 = 3{,}04$). La tasa de error sigue siendo
la variable dominante, lo que respalda el razonamiento de la página 6, pero el
modelo exagera esa dominancia en más del doble.

\subsection*{03. Unidades y rango de validez}

Hay dos problemas adicionales, y son más graves que el anterior.

El símbolo $p$ cambia de unidad dentro de la misma página. En la matriz
morfológica está en por ciento, con rangos de 0,003 a 0,3, que se corresponden
con las fidelidades declaradas. En el análisis de sensibilidad está expresado
como fracción, con valores de 0,001 a 0,1, y el punto nominal se enuncia como
$p = 5$ por ciento. Es decir, el punto de diseño nominal tiene una tasa de
error entre 17 y 1250 veces peor que cualquiera de los tres dispositivos de la
matriz. El número de cúbits nominal, $n = 128$, es cuatro veces mayor que la
opción más grande de la matriz. El punto de diseño sobre el que se calculan
las sensibilidades queda fuera del espacio de diseño que la propia página
define.

El segundo problema es el rango de validez del modelo. El Cuadro~\ref{tab:qv}
compara el volumen cuántico que predice la expresión $2^{n}(1-p)^{2}$ con el
que la propia hoja de ruta reporta para cada máquina.

\begin{table}[ht]
  \centering
  \small
  \caption{Volumen cuántico predicho por el modelo de la hoja de ruta frente al
  volumen cuántico que ella misma reporta. Cálculo propio a partir de los datos
  de \textcite{numa2023}, tomando $p$ igual a la tasa de error de puerta de dos
  cúbits.}
  \label{tab:qv}
  \begin{tabular}{@{}lrrrr@{}}
    \toprule
    \textbf{Dispositivo} & \textbf{$n$} & \textbf{$p$} &
    \textbf{$2^{n}(1-p)^{2}$} & \textbf{QV reportado} \\
    \midrule
    H1-2 (Quantinuum)   &  12 & 0,003 & $4{,}07\times10^{3}$  & $4\,096$ \\
    H1-1 (Quantinuum)   &  20 & 0,002 & $1{,}04\times10^{6}$  & $524\,288$ \\
    H2 (Quantinuum)     &  32 & 0,002 & $4{,}28\times10^{9}$  & $65\,536$ \\
    Aspen-M-1 (Rigetti) &  80 & 0,054 & $1{,}08\times10^{24}$ & $8$ \\
    \bottomrule
  \end{tabular}
\end{table}

El modelo acierta con el H1-2, se equivoca por un factor de dos con el H1-1, se
equivoca por un factor de sesenta y cinco mil con el H2 y se equivoca por
veintitrés órdenes de magnitud con el Aspen-M-1. El patrón es claro: la
expresión sólo se aproxima al comportamiento observado mientras el número de
cúbits es pequeño, porque no incorpora la limitación de profundidad de circuito
que en la práctica acota el volumen cuántico. Dicho de otro modo, el modelo
funciona para el dispositivo que se quiere elegir y falla para los demás, y el
punto nominal de $n = 128$ está muy lejos de su zona de validez.

Esto tiene una consecuencia directa sobre la pregunta que motiva este
comentario: el modelo técnico no puede usarse para proyectar nada. Es una
fotografía estática, sin variable temporal, cuya extrapolación produce valores
absurdos.

\section{El modelo financiero (elemento 8)}

\subsection*{01. Estructura y consistencia aritmética}

Conviene empezar por lo positivo, porque es notable. El modelo financiero es
íntegramente reproducible. Consta de un análisis de cuatro compañías
comparables (NVIDIA, IBM, Microsoft y Google), una proyección anual de catorce
periodos que abarca de 2022 a 2035, y un descuento de flujos con tasa del 10
por ciento, impuesto del 21 por ciento y crecimiento terminal del 5 por ciento.

Se recalcularon todas las líneas. El tamaño de mercado, la cuota, los ingresos,
el valor terminal calculado como $CF_{14}/(r-g) = 4\,658/0{,}05 = 93\,160$, su
valor presente de $26\,985$ y el valor presente neto total de $32\,210$
coinciden con las cifras publicadas de $93\,164$, $26\,986$ y $32\,211$ dentro
del redondeo. La tabla de comparables también es internamente consistente: los
ingresos de aprendizaje automático son el producto de las ventas por el
porcentaje declarado en cada caso.

Merece señalarse un acierto metodológico que a primera vista parece un error:
el flujo del periodo 14 no se descuenta ni se suma, y eso es correcto, porque
el valor terminal situado en el periodo 13 ya capitaliza ese flujo y los
siguientes. La aritmética del descuento está bien planteada.

\subsection*{02. El año base y la cuota de mercado}

El problema del modelo no está en su aritmética sino en sus entradas, y es
grave. El Cuadro~\ref{tab:financiero} extrae las líneas decisivas.

\begin{table}[ht]
  \centering
  \small
  \caption{Líneas seleccionadas del modelo financiero, en millones de dólares.
  Transcripción de \textcite{numa2023}.}
  \label{tab:financiero}
  \begin{tabular}{@{}lrrrrr@{}}
    \toprule
    & \textbf{2022} & \textbf{2025} & \textbf{2030} & \textbf{2032} &
    \textbf{2035} \\
    \midrule
    Tamaño de mercado & 613,0 & 1\,346,8 & 5\,000,4 & 7\,800,7 & 12\,917,9 \\
    Cuota de mercado  & 10,1\% & 13,4\% & 21,6\% & 26,1\% & 34,8\% \\
    Ingresos          & 62 & 181 & 1\,079 & 2\,037 & 4\,491 \\
    \bottomrule
  \end{tabular}
\end{table}

El modelo asigna a la empresa una cuota del 10,1 por ciento del mercado de
aprendizaje automático cuántico en 2022 y unos ingresos de 62 millones de
dólares en ese mismo año. La nota de supuestos indica que ese 10,1 por ciento
procede del \enquote{promedio de la industria}, y en efecto coincide con el
10,07 por ciento que la tabla de comparables calcula como cuota media de
NVIDIA, IBM, Microsoft y Google.

Ese encadenamiento no se sostiene, por tres razones acumulativas.

Primera, la hoja de ruta describe en la página 6 una empresa que
\enquote{pretende lanzar} el producto. El modelo financiero le atribuye
ingresos desde el primer periodo, sin año de lanzamiento, sin rampa y sin
ningún ejercicio con ingresos nulos.

Segunda, la cuota inicial que se le asigna es diez veces el objetivo
estratégico que la propia hoja de ruta se fija, más del 1 por ciento, y que
califica de \enquote{en riesgo}. La cuota final, 34,8 por ciento, es treinta y
cinco veces ese objetivo y supera a la del comparable más alto de su propia
tabla, que es Google con el 28,21 por ciento.

Tercera, los denominadores no son los mismos. Si se dividen los ingresos de
aprendizaje automático de cada comparable por su cuota declarada, resulta de
forma consistente un mercado total de unos 158 mil millones de dólares, que es
el mercado de aprendizaje automático clásico. La cuota que el modelo aplica es
en cambio del mercado de aprendizaje automático cuántico, que en la misma tabla
vale 613 millones en 2022. Se traslada un porcentaje entre dos mercados que
difieren en más de dos órdenes de magnitud.

Además, la cuota crece de forma mecánica al 10 por ciento anual durante trece
periodos, sin ningún vínculo con los hitos técnicos de 2025, 2030 o 2032. Los
ingresos suben igual se cumplan o no. Este es el hallazgo decisivo para la
pregunta de la proyección temporal y se retoma más adelante.

\subsection*{03. El signo del gasto de capital}

El modelo fija el gasto de capital en el 92,63 por ciento de las ventas, valor
tomado del promedio de la tabla de comparables, y lo \emph{suma} al flujo de
caja. La fórmula que la propia hoja de ruta enuncia en su columna de supuestos
es
\[
  CF = (1-\tau)\,EBITDA + \tau\,\text{Depreciación} - CapEx -
  \Delta\text{Capital de trabajo},
\]
donde el gasto de capital se resta. Aplicando el signo tal como la fórmula lo
declara, y con las cifras de la propia tabla, el flujo de 2022 sería
$18 + 1 - 57 - 4 = -42$ millones, el de 2035 sería $-3\,662$ millones y todos
los años intermedios serían negativos, porque un gasto de capital del 92,6 por
ciento de las ventas excede con holgura un EBITDA después de impuestos que
ronda el 28,5 por ciento de las ventas. El valor presente neto resultante sería
de unos $-25\,236$ millones en lugar de $+32\,211$. El signo de esa línea
decide por sí solo si el proyecto vale treinta y dos mil millones o pierde
veinticinco mil.

El origen del problema está aguas arriba, en la tabla de comparables, donde el
gasto de capital alcanza el 119,2 por ciento de las ventas para IBM y el 119,6
por ciento para Microsoft. Ninguna compañía madura sostiene una inversión anual
en activos superior a sus ingresos, de modo que esa columna no puede ser gasto
de capital anual, y el promedio del 92,63 por ciento que de ella se deriva no
es un parámetro utilizable.

En la misma tabla es visible un error de hoja de cálculo sin resolver,
\enquote{\#REF!}, en la fila de depreciación sobre activos fijos.

\subsection*{04. La concentración en el valor terminal}

De los 32\,211 millones de valor presente neto, 26\,986 son valor terminal, es
decir el 83,8 por ciento. Los flujos explícitos de los trece años proyectados
aportan sólo 5\,225 millones. La valoración no descansa sobre el periodo que la
hoja de ruta modela, sino sobre una perpetuidad que empieza donde el modelo
termina.

Para el objeto de este comentario, esto significa que el modelo financiero no
es una proyección de cuándo la tecnología generará valor, sino una apuesta
sobre las condiciones posteriores a 2035. El horizonte 2035 no es una fecha
derivada: es el punto donde se corta el análisis.

Hay finalmente dos discrepancias menores de atribución y de narrativa. La
tabla acredita el tamaño de mercado a Statista, mientras la viñeta de la misma
página lo acredita a Virtue Market Research; las cifras (613,0 millones en 2022
y 5\,000,4 millones en 2030) coinciden exactamente con el informe de Virtue
Market Research, cuyo horizonte publicado termina en 2030
\parencite{virtue2023}, de manera que las tasas aplicadas de 2031 a 2035 no
proceden de la fuente citada. Y la viñeta afirma un crecimiento del 30 por
ciento hasta 2030 que se moderaría al 20 por ciento hasta 2035, mientras la
tabla aplica en realidad el 30 por ciento hasta 2031, el 20 por ciento entre
2032 y 2034 y el 15 por ciento en 2035.

\section{Portfolio de I+D, publicaciones y patentes (elementos 9 y 10)}

\subsection*{01. El portfolio ausente}

La página 7 se titula \enquote{R\&D Projects, Company Positioning vs.
Competition with FOMs}, pero no contiene ningún proyecto de la empresa.
Contiene hardware de otros fabricantes y una ilustración generativa. No hay
lista de proyectos propios, ni asignación de presupuesto, ni calendario, ni
responsables, ni valoración por proyecto, ni matriz de riesgo frente a retorno,
ni prototipos.

El elemento 9 es, junto al 12, el que más pesa en la pregunta temporal. Un
portfolio de proyectos con fechas de inicio y fin es el instrumento que traduce
una intención estratégica en un calendario verificable. Sin él, los hitos de la
página 4 quedan sin nada debajo que los sostenga.

\subsection*{02. Publicaciones y patentes}

La genealogía de publicaciones sí es un buen elemento. Recoge once hitos
fechados entre 1995 y 2021: los primeros trabajos de inspiración biológica
sobre aprendizaje automático cuántico (Kak, 1995), la primera propuesta de red
neuronal cuántica competitiva (Ventura, 1999), la neurona tipo cúbit (Matsui y
otros, 2000), el perceptrón cuántico (Altaiski, 2001), las máquinas de vectores
de soporte cuánticas (Anguita, 2003), el resultado sobre patrones que las redes
clásicas no pueden aprender (Liu, 2013), las $K$-medias cuánticas (Lloyd,
2013), la primera monografía sobre la materia (Wittek, 2014), el modelo de red
neuronal cuántica para procesadores de corto plazo (Farhi, 2017), la primera
implementación de una neurona cuántica real en un procesador (Tacchino y otros,
2019) y la primera versión de TensorFlow Quantum (Broughton y otros, 2020). La
hoja de ruta divide con acierto la evolución en dos fases, formulación de
modelos desde mediados de los años noventa hasta 2007 y énfasis posterior en la
implementación, siguiendo a \textcite{alchieri2021}.

Sobre las dos patentes hay una corrección que hacer. La hoja de ruta fecha
ambas en 2022. La patente de Google, US11449760B2, \enquote{Quantum assisted
optimization}, fue efectivamente concedida el 20 de septiembre de 2022, a
partir de una solicitud con prioridad del 25 de abril de 2016
\parencite{patentGoogle2022}. La patente de Accenture, US10275721B2,
\enquote{Quantum computing machine learning module}, fue concedida el 30 de
abril de 2019, no en 2022, a partir de una solicitud del 19 de abril de 2017
\parencite{patentAccenture2019}.

La corrección no es un detalle bibliográfico. La hoja de ruta usa años de
concesión como marcadores de progreso, y los años de concesión son un
indicador ruidoso porque dependen de la tramitación: en estos dos casos median
dos y seis años entre solicitud y concesión. Si se quiere leer la actividad de
patentes como señal temporal, la fecha pertinente es la de prioridad.

\section{Estrategia tecnológica y madurez (elementos 11 y 12)}

\subsection*{01. La declaración de estrategia}

La página 4 es el único lugar de toda la hoja de ruta donde aparecen fechas
objetivo. El enunciado general fija avances computacionales de gran alcance
para 2035, y las tres líneas estratégicas y el gráfico que las acompaña
producen la escalera de hitos que recoge la Figura~\ref{fig:hitos}.

\begin{figure}[ht]
  \centering
  \begin{tikzpicture}[
    hito/.style={draw, rounded corners, align=center, font=\scriptsize,
      text width=0.18\textwidth, inner sep=4pt, minimum height=13mm}
  ]
    \draw[-{Stealth[length=3mm]}, very thick, gray!60]
      (0.01\textwidth,0) -- (0.96\textwidth,0);
    \foreach \x/\a in {0.12/2023, 0.2467/2025, 0.5633/2030, 0.69/2032, 0.88/2035}
      {\draw[thick] (\x\textwidth,-1.2mm) -- (\x\textwidth,1.2mm);
       \node[above=1mm, font=\scriptsize\bfseries] at (\x\textwidth,1.5mm) {\a};}
    \node[hito, below=9mm, fill=gray!8] at (0.12\textwidth,0)
      {Hoja de ruta\\redactada};
    \node[hito, above=6mm, fill=blue!7] at (0.2467\textwidth,0)
      {Plataforma de\\computación cuántica\\robusta para IA/ML};
    \node[hito, above=6mm, fill=blue!7] at (0.5633\textwidth,0)
      {Algoritmos cuánticos\\que superan a los\\clásicos};
    \node[hito, below=9mm, fill=orange!12] at (0.69\textwidth,0)
      {Sistemas híbridos\\cuántico-clásicos\\\emph{operativos}};
    \node[hito, above=6mm, fill=green!10] at (0.88\textwidth,0)
      {Capacidades\\computacionales\\transformadoras};
    \draw[-{Stealth[length=2mm]}, thick, orange!70!black]
      (0.6\textwidth,-9mm) .. controls (0.5633\textwidth,-6mm) ..
      (0.5633\textwidth,-1.5mm);
  \end{tikzpicture}
  \caption{Escalera de hitos reconstruida a partir de la declaración de
  estrategia tecnológica. El paso descrito como \emph{solución intermedia}
  hacia los sistemas plenamente cuánticos se fecha en 2032, después del hito de
  2030 al que debería conducir. Elaboración propia a partir de
  \textcite{numa2023}.}
  \label{fig:hitos}
\end{figure}

La declaración tiene dos defectos internos.

El primero es de secuencia. La tercera línea estratégica describe los sistemas
híbridos cuántico-clásicos como \enquote{una solución intermedia para los
avances en IA/ML y un paso hacia soluciones plenamente cuánticas}. El gráfico
los fecha en 2032, es decir, dos años \emph{después} del hito de 2030 en el que
los algoritmos cuánticos ya habrían superado a los clásicos. Un paso intermedio
que llega después del destino al que conduce indica que la escalera no se
construyó como una secuencia de dependencias sino como una lista de
aspiraciones ordenadas por ambición.

El segundo es de consistencia entre soporte y texto. La fecha de 2032 aparece
únicamente en el gráfico; la viñeta correspondiente no lleva fecha. Las de 2025
y 2030 sí aparecen en ambos.

Y hay un defecto de método que afecta a las cuatro fechas: ninguna se deriva de
nada. No se declara qué valor de qué figura de mérito define el cumplimiento de
cada hito, ni con qué tasa de mejora se llegaría a él, ni de qué proyecto
depende. Son fechas afirmadas, no calculadas.

\subsection*{02. La evaluación de madurez ausente}

El elemento 12 no está. No hay evaluación del nivel de madurez de la propia
hoja de ruta ni de sus subsistemas, ni una escala de tipo TRL, ni una
autoevaluación de la completitud de los elementos anteriores.

Esta ausencia es la más costosa de todas para la pregunta temporal, porque un
nivel de madurez es exactamente la magnitud que se convierte en tiempo hasta la
disponibilidad. Sin saber en qué nivel está cada componente, no hay origen
desde el que contar.

\section{¿Hay material para proyectar cuándo estará disponible la tecnología?}

Esta es la pregunta que motiva el comentario y merece una respuesta ordenada.
La respuesta corta es que hay material fechado, pero no material proyectable, y
que la brecha entre ambas cosas es identificable con precisión.

\subsection*{01. Inventario del material fechado}

La hoja de ruta contiene cinco conjuntos de datos con eje temporal:

\begin{enumerate}
  \item \textbf{Los hitos estratégicos} de 2025, 2030, 2032 y 2035
    (Figura~\ref{fig:hitos}).
  \item \textbf{La serie financiera anual} de 2022 a 2035, con tamaño de
    mercado, cuota, ingresos y flujos de caja para cada ejercicio.
  \item \textbf{La tabla competitiva}, con año de lanzamiento entre 2017 y 2023
    para veinte procesadores, cúbits físicos para todos ellos y volumen
    cuántico para cuatro.
  \item \textbf{La genealogía de publicaciones}, con once hitos entre 1995 y
    2021, más dos patentes.
  \item \textbf{La curva de cómputo de entrenamiento}, con cómputo en FLOPS
    frente a fecha de publicación entre 2010 y 2022.
\end{enumerate}

Es más material del que suele encontrarse en un ejercicio de este tipo. El
problema es lo que falta para convertirlo en una proyección.

\subsection*{02. Los cuatro ingredientes que faltan}

Una proyección de disponibilidad requiere cuatro cosas: un umbral, una línea
base, una tasa de mejora y un vínculo entre el umbral técnico y la
disponibilidad comercial. La hoja de ruta no tiene ninguna de las cuatro.

\textbf{No hay umbral.} Ninguna de las nueve figuras de mérito tiene valor
objetivo. \enquote{Disponible} no está definido en términos medibles. No se
puede fechar el cruce de una línea que no se ha dibujado.

\textbf{No hay línea base para lo que importa.} El único dato de capacidad con
eje de calendario es de hardware, cúbits y volumen cuántico, y no de
rendimiento de aprendizaje automático cuántico. El único dato de rendimiento de
aprendizaje automático cuántico, el gráfico de pérdida frente a iteraciones
tomado de \textcite{huang2021}, tiene en el eje horizontal iteraciones de
entrenamiento y no años. Ninguno de los tres gráficos de la página 3 tiene
simultáneamente una figura de mérito de la tecnología en el eje vertical y
tiempo de calendario en el horizontal.

\textbf{No hay tasa.} En ninguna página se estima una velocidad de mejora. El
modelo técnico calcula derivadas respecto a $n$ y a $p$, pero no respecto al
tiempo. No hay ningún $dFOM/dt$ en todo el documento.

\textbf{No hay vínculo.} El modelo financiero, que sí tiene calendario anual
completo, es independiente de los hitos técnicos. Los ingresos de 2031 son los
mismos si la plataforma de 2025 se retrasa cinco años. Esto es lo que convierte
la serie financiera en un calendario y no en una proyección: nada en ella está
condicionado a la madurez de la tecnología.

\subsection*{03. Lo que sí permiten los datos de la propia hoja de ruta}

Pese a lo anterior, la tabla competitiva contiene la única serie de la hoja de
ruta que empareja una capacidad técnica con años de calendario, y sí admite un
ajuste. Los doce procesadores de Rigetti que la tabla recoge cubren de 2017 a
2022 con recuento de cúbits y fidelidad de puerta de dos cúbits. Tomando el
mejor valor de cada año y ajustando una tendencia exponencial se obtienen los
resultados de la Figura~\ref{fig:tendencias} y del Cuadro~\ref{tab:tasas}.

\begin{figure}[ht]
  \centering
  \begin{tikzpicture}
    \begin{semilogyaxis}[
      width=0.47\textwidth, height=5.1cm,
      xlabel={Año de lanzamiento}, ylabel={Cúbits físicos},
      xmin=2016.5, xmax=2022.5, ymin=10, ymax=120,
      xtick={2017,2018,2019,2020,2021,2022},
      xticklabel style={/pgf/number format/1000 sep=},
      log basis y=10, ytick={10,20,50,100},
      yticklabels={10,20,50,100},
      grid=both, grid style={gray!18}, tick label style={font=\scriptsize},
      label style={font=\scriptsize}, title style={font=\scriptsize\bfseries},
      title={Recuento de cúbits: duplica cada 2,4 años}]
      \addplot[only marks, mark=*, mark size=1.6pt, blue!70!black]
        coordinates {(2017,19) (2018,16) (2019,28) (2020,31) (2021,40) (2022,80)};
      \addplot[thick, dashed, red!70!black] coordinates {(2017,15.0) (2022,63.0)};
    \end{semilogyaxis}
  \end{tikzpicture}\hfill
  \begin{tikzpicture}
    \begin{axis}[
      width=0.47\textwidth, height=5.1cm,
      xlabel={Año de lanzamiento},
      ylabel={Error de puerta de 2 cúbits [\%]},
      xmin=2016.5, xmax=2022.5, ymin=0, ymax=14,
      xtick={2017,2018,2019,2020,2021,2022},
      xticklabel style={/pgf/number format/1000 sep=},
      grid=both, grid style={gray!18}, tick label style={font=\scriptsize},
      label style={font=\scriptsize}, title style={font=\scriptsize\bfseries},
      title={Tasa de error: sin mejora desde 2019}]
      \addplot[only marks, mark=*, mark size=1.6pt, blue!70!black]
        coordinates {(2017,12.5) (2018,9.16) (2019,4.8) (2020,5.66)
                     (2021,5.34) (2022,4.9)};
      \addplot[thick, dashed, red!70!black]
        coordinates {(2017,10.20) (2022,4.25)};
      \addplot[thick, orange!85!black] coordinates {(2019,5.175) (2022,5.175)};
    \end{axis}
  \end{tikzpicture}
  \caption{Tendencias derivadas de la serie de doce procesadores de Rigetti que
  la hoja de ruta recoge, tomando el mejor valor de cada año. La línea
  discontinua es el ajuste exponencial 2017--2022; la línea naranja continua es
  la media del subperiodo 2019--2022. Cálculo propio a partir de los datos de
  \textcite{numa2023}.}
  \label{fig:tendencias}
\end{figure}

\begin{table}[ht]
  \centering
  \small
  \caption{Tasas de mejora ajustadas sobre los datos de la hoja de ruta.
  Cálculo propio.}
  \label{tab:tasas}
  \begin{tabularx}{\textwidth}{@{}>{\raggedright\arraybackslash}X rrr@{}}
    \toprule
    \textbf{Serie} & \textbf{Tasa anual} & \textbf{Tiempo de duplicación
    o reducción a la mitad} & \textbf{$R^{2}$} \\
    \midrule
    Cúbits físicos, 2017--2022 & $+33{,}2\%$ & 2,42 años & 0,872 \\
    Error de dos cúbits, 2017--2022 & $-16{,}1\%$ & 3,95 años & 0,695 \\
    Error de dos cúbits, 2019--2022 & $+0{,}04\%$ & no aplicable & 0,000 \\
    \bottomrule
  \end{tabularx}
\end{table}

De aquí se desprenden tres conclusiones, y las tres apuntan en contra de la
posibilidad de fechar la disponibilidad con este material.

\textbf{Primera: las dos variables del modelo técnico avanzan a velocidades
incomparables.} El recuento de cúbits crece de forma robusta, duplicándose cada
2,4 años con un ajuste razonable. La tasa de error, en cambio, mejoró con
fuerza entre 2017 y 2019 y desde entonces no se mueve: el subperiodo 2019--2022
tiene una pendiente indistinguible de cero y un $R^{2}$ nulo. Y la tasa de
error es, según la propia hoja de ruta, \enquote{el factor más significativo}
al evaluar la eficiencia del hardware, además de la variable a la que su
análisis de sensibilidad atribuye entre tres y siete veces más apalancamiento.
La hoja de ruta apuesta el rendimiento del sistema a la variable que, en sus
propios datos, no está mejorando.

\textbf{Segunda: extrapolar hacia el objetivo estratégico da fechas
incompatibles con los hitos.} Con la tendencia 2017--2022, pasar del 4,9 por
ciento de error del mejor dispositivo de 2022 al 0,1 por ciento que exige el
objetivo de fidelidad superior al 99,9 por ciento requeriría 22,2 años, es
decir, alrededor de 2044. Con la tendencia 2019--2022, no se llegaría nunca. La
hoja de ruta sitúa la plataforma robusta en 2025.

\textbf{Tercera: el modelo técnico se rompe si se le aplica la única tendencia
disponible.} Si se proyecta el recuento de cúbits con su tasa observada y se
mantiene la tasa de error en su valor estancado, la expresión
$QV = 2^{n}(1-p)^{2}$ predice para 2024 un volumen cuántico del orden de
$2^{160}$, cuando la propia tabla de la hoja de ruta atribuye al procesador de
80 cúbits de 2022 un volumen cuántico de 8. La extrapolación del modelo técnico
sobre datos del propio documento produce un resultado absurdo en dos años.

Conviene añadir una cautela que refuerza el argumento en lugar de debilitarlo.
Estas tendencias se han ajustado sobre una sola familia de dispositivos de un
solo fabricante y una sola arquitectura, la superconductora, porque es la única
serie con suficientes puntos en la tabla. La comparación con la otra
arquitectura presente lo confirma: el H1-1 de iones atrapados alcanzaba en 2022
una fidelidad de dos cúbits del 99,8 por ciento, muy por delante de esa
tendencia. Es decir, el salto de rendimiento relevante no llegó por avance
temporal dentro de una arquitectura, sino por cambio de arquitectura. Esto es
precisamente lo que invalida la extrapolación temporal como instrumento: el
progreso en este campo es discontinuo y depende de qué tecnología física gane,
una cuestión que ninguna curva de tendencia responde. La tabla competitiva,
además, sólo tiene volumen cuántico en cuatro de veinte filas, todas de 2022 y
2023, y con valores no monótonos, ya que el H1-1 de 2022 declara 524\,288
frente a los 65\,536 del H2 de 2023. Con dos años y valores decrecientes no hay
ninguna tendencia que ajustar.

\subsection*{04. Por qué la restricción probablemente no es de hardware}

Hay un argumento más profundo, y procede de una de las fuentes que la hoja de
ruta cita.

La página 3 organiza su razonamiento en tres pasos: el aprendizaje automático
necesita procesar datos a escala cada vez mayor; la evolución de la computación
cuántica mejorará su rendimiento; la evolución de los algoritmos cuánticos
mejorará su eficiencia. El segundo paso se apoya en un gráfico de tiempo de
ejecución frente a tamaño del problema, atribuido a \enquote{Florian Meyer, ETH
Zurich} en \emph{Communications of the ACM}, 2023.

La verificación de esa referencia arroja dos resultados. El primero es de
autoría: el artículo de \emph{Communications of the ACM} de mayo de 2023 se
titula \enquote{Disentangling Hype from Practicality: On Realistically
Achieving Quantum Advantage} y es obra de Torsten Hoefler, Thomas Häner y
Matthias Troyer \parencite{hoefler2023}. Florian Meyer es el autor de la nota
divulgativa que ETH Zurich publicó sobre ese artículo, no del artículo.

El segundo resultado es de fondo y es el que importa. La tesis de ese trabajo
es que la ventaja cuántica práctica exige problemas de \emph{datos pequeños} y
aceleraciones \emph{superiores a las cuadráticas}, que la mayoría de los
algoritmos y aplicaciones cuánticas propuestos no alcanzan esas aceleraciones,
y que el potencial se concentra por ahora en ciencia de materiales y química
\parencite{hoefler2023}. La hoja de ruta invoca ese trabajo para sostener que
la era del aprendizaje automático cuántico se acerca, cuando el trabajo citado
sostiene que el requisito es de datos pequeños, justo lo contrario de la
premisa \enquote{el aprendizaje automático necesita procesar datos a gran
escala} con la que abre la misma página. La fuente citada contradice el
argumento para el que se la cita.

Si el trabajo citado tiene razón, la consecuencia metodológica es directa: la
restricción determinante no es el hardware sino la existencia de algoritmos con
la clase de aceleración adecuada sobre problemas del tamaño adecuado. Y la
llegada de un resultado algorítmico no es extrapolable a partir de tendencias
de cúbits ni de ninguna otra curva de capacidad física. Este es el motivo más
sólido para no ofrecer una fecha: no es que falten datos de hardware, es que
la variable crítica es de naturaleza discreta e impredecible.

Merece señalarse, en el mismo sentido, que el gráfico del propio artículo no
tiene eje de calendario. Su eje horizontal es el tamaño del problema. Para
convertir ese cruce en una fecha haría falta conocer cómo crece con el tiempo
el tamaño de los problemas de los clientes, algo que la hoja de ruta no aporta.
Es notable que el documento tenga la curva de demanda, el cómputo de
entrenamiento de los modelos clásicos frente a la fecha de publicación, y tenga
el concepto de cruce, pero no llegue a dibujar la curva de oferta con eje
temporal que permitiría intersecarlas. Se queda a un paso de la construcción
que respondería a la pregunta.

\subsection*{05. Veredicto}

La hoja de ruta contiene una \emph{declaración de intención} fechada en 2035 y
una escalera de hitos, pero no contiene el material para una proyección
independiente de disponibilidad. La distinción es la siguiente: una declaración
de intención afirma una fecha; una proyección la deriva de una línea base, una
tasa y un umbral, y expone el intervalo de incertidumbre que la acompaña.

Con el contenido del documento, el enunciado más fuerte que se puede sostener
con rigor es direccional y no está fechado: sobre los datos que la propia hoja
de ruta reúne, el recuento de cúbits mejora de forma sostenida mientras la
tasa de error, que su análisis identifica como la variable dominante, no ha
mejorado desde 2019; y la fuente que el documento invoca para justificar el
optimismo sostiene que el cuello de botella es algorítmico. Nada de esto
permite escribir un año, y presentarlo como si lo permitiera sería precisamente
el error que el trabajo citado advierte.

Los hitos de 2025, 2030 y 2032 pueden evaluarse en retrospectiva desde la
posición actual, y esa es una vía legítima para valorar la calidad de la hoja
de ruta como ejercicio de previsión, pero no es una proyección hecha a partir
de ella.

\section{Recomendaciones}

Las siguientes seis medidas cerrarían la brecha entre el material fechado que
la hoja de ruta ya tiene y una proyección defendible. Están ordenadas por
impacto sobre la pregunta temporal.

\begin{enumerate}[label=\textbf{\arabic*.}]
  \item \textbf{Definir la disponibilidad de forma operativa.} Una figura de
    mérito, un umbral y una carga de trabajo. Por ejemplo: exactitud de modelo
    igual o superior a un valor determinado sobre una tarea concreta de un
    cliente, con un coste total por inferencia por debajo de un límite, frente
    a la mejor alternativa clásica. Sin esta frase, ningún elemento posterior
    puede fechar nada.
  \item \textbf{Reconciliar los tres vocabularios de métricas.} Incorporar el
    volumen cuántico, la fidelidad de dos cúbits y el recuento de cúbits a la
    lista de figuras de mérito, o bien retirarlos del modelo técnico, y
    declarar la relación entre cada objetivo estratégico y la figura de mérito
    que lo mide. Conviene además reducir las nueve figuras a las tres o cuatro
    que realmente se van a seguir, con valor de partida y valor objetivo
    fechado para cada una.
  \item \textbf{Construir la curva de oferta con eje temporal.} La misma fuente
    que la hoja de ruta ya utiliza permite completar el volumen cuántico por
    año para los principales fabricantes, e incorporar a IBM, Google, Intel y
    D-Wave, hoy ausentes del extracto. A ello pueden añadirse las hojas de ruta
    publicadas por los propios fabricantes, que sí están fechadas: la de IBM de
    septiembre de 2020 fijaba Hummingbird con 65 cúbits en 2020, Eagle con 127
    en 2021, Osprey con 433 en 2022 y Condor con 1121 en 2023
    \parencite{ibm2020}, y el 4 de diciembre de 2023, un día antes de la fecha
    de creación del archivo analizado, se amplió hasta 2033 con Starling
    previsto para 2029 \parencite{ibm2023}. Superponer esa curva de oferta a la
    curva de demanda de cómputo que la hoja de ruta ya incluye permite
    localizar un cruce y acompañarlo de una banda de incertidumbre.
  \item \textbf{Añadir los dos elementos ausentes.} Un portfolio de proyectos
    propios con fechas de inicio y fin, presupuesto y dependencias, y una
    evaluación de madurez por subsistema. Son los dos elementos que convierten
    una aspiración en un calendario auditable, y son precisamente los dos que
    faltan.
  \item \textbf{Condicionar el modelo financiero a los hitos técnicos.} Fijar
    un año de lanzamiento, empezar la serie con ingresos nulos, partir de una
    cuota coherente con el objetivo declarado de más del 1 por ciento en lugar
    del 10,1 por ciento heredado de los comparables, corregir el signo del
    gasto de capital y ejecutar al menos un escenario con los hitos retrasados.
    Dado que el 83,8 por ciento del valor presente neto reside en el valor
    terminal, un análisis de sensibilidad sobre la fecha de disponibilidad es
    más informativo que cualquier refinamiento de los flujos explícitos.
  \item \textbf{Releer las fuentes citadas.} Corregir la atribución del
    artículo de \emph{Communications of the ACM}, sustituir la entrada
    divulgativa que respalda la curva de cómputo por el estudio primario
    \parencite{sevilla2022}, y ajustar la fecha de la patente de Accenture. Más
    importante que las correcciones: si la tesis de \textcite{hoefler2023} se
    acepta, la hoja de ruta debería reorientar su premisa desde el
    procesamiento de datos a gran escala hacia problemas de datos pequeños con
    aceleraciones superiores a las cuadráticas, lo que cambiaría tanto el
    portfolio de I+D como el mercado objetivo.
\end{enumerate}

\section{Conclusiones}

La hoja de ruta analizada es un ejercicio sólido en su mitad estratégica y
frágil en su mitad predictiva. La correspondencia con las motivaciones de la
empresa está bien construida, con objetivos verificables y un semáforo que no
oculta los incumplimientos; la genealogía de publicaciones es útil; y el modelo
financiero, cuya aritmética se reproduce sin discrepancias, revela una
disciplina de cálculo apreciable. Nueve de los doce elementos están presentes
de forma sustantiva.

El defecto que atraviesa el documento es la falta de trazabilidad entre
métricas. La estrategia se evalúa con cuatro objetivos, el sistema se describe
con nueve figuras de mérito y la decisión técnica se toma con el volumen
cuántico, sin que ninguno de los tres conjuntos se relacione con los otros.
Cuando esa cadena falta, cada elemento puede parecer correcto por separado y el
conjunto no puede responder a la pregunta para la que existe una hoja de ruta.

Respecto a la pregunta planteada, la conclusión es que hay material fechado
pero no material proyectable. Faltan los cuatro ingredientes de una proyección:
un umbral que defina la disponibilidad, una línea base de la magnitud
pertinente, una tasa de mejora y un vínculo entre el umbral técnico y el
calendario comercial. Y faltan, además, los dos elementos que habrían aportado
calendario: el portfolio de proyectos y la evaluación de madurez.

La ironía del caso es que la propia hoja de ruta contiene los datos con los que
demostrarlo. Su tabla competitiva permite establecer que el recuento de cúbits
se duplica cada 2,4 años mientras la tasa de error, la variable que su análisis
de sensibilidad señala como dominante, no ha mejorado desde 2019; su modelo
técnico produce resultados absurdos en cuanto se le aplica esa tendencia; y la
fuente que invoca para sostener que la era del aprendizaje automático cuántico
se acerca sostiene en realidad que la mayoría de los algoritmos cuánticos
propuestos no alcanzan la aceleración necesaria para ser prácticos, y que el
requisito es de datos pequeños, lo contrario de la premisa con la que el
documento abre.

De ahí se sigue la lección de gestión tecnológica que este caso ofrece. La
tentación de una hoja de ruta es fijar la fecha primero y buscar después el
respaldo. El orden correcto es el inverso: definir el umbral, medir la línea
base, estimar la tasa y sólo entonces derivar la fecha, con la incertidumbre
declarada. Una hoja de ruta que no puede fechar su objetivo sigue siendo útil
si dice con precisión qué le falta para poder hacerlo. Ese es, en definitiva,
el aporte principal de este comentario sobre el documento analizado.

\printbibliography[title={Referencias bibliográficas}]

\end{document}
